problem:  
Let \((X,d,\mu)\) be a complete, doubling metric measure space supporting a \(p\)-Poincaré inequality for some \(p>1\), so that \(X\) is a Lipschitz differentiability space in the sense of Cheeger.  
Suppose that for \(\mu\)-a.e. \(x\in X\), every tangent cone \(\operatorname{Tan}(X,d,\mu,x)\) is a Carnot group \((G_x,d_x,\mu_x)\) equipped with its Haar measure.  

**Open problem:**  
Can there exist a doubling PI space \((X,d,\mu)\) for which the Lie algebra type of \(G_x\) **varies** with \(x\)?  
Equivalently, is there a metric-measure space whose pmGH tangents are Carnot groups at almost every point, yet the associated graded Lie algebras \(\mathfrak g_x = \operatorname{Lie}(G_x)\) are not all isomorphic?  
If such examples exist, can one construct them explicitly—perhaps as measured Gromov–Hausdorff limits of degenerating families of equiregular sub-Riemannian manifolds with varying horizontal distributions?  
If they do not exist, can one prove a **rigidity theorem** stating that the Carnot type must be locally constant whenever tangents are Carnot groups a.e.?

Why is it a "good" problem:  
This formulation corrects the feasibility concerns: before asking how to glue variable tangents, it first poses the fundamental existence question—does variation of Carnot tangent type even occur?  
This question stands at the interface of analysis on metric spaces, sub-Riemannian geometry, and group-theoretic classification of nilpotent Lie algebras. It pushes against what is currently known: in all established examples (Riemannian manifolds, equiregular sub-Riemannian manifolds, Carnot-group limits of such), the infinitesimal Carnot type is locally constant. Verifying whether this is enforced by the PI and doubling structure, or merely a feature of known constructions, would illuminate how rigid the sub-Riemannian tangent behavior truly is.  

A positive answer (existence of varying tangent types) would reveal an entirely new class of metric measure geometries—“non-equiregular PI spaces”—suggesting the need for a new analytic framework.  
A negative answer would establish a strong **homogeneity rigidity principle** for PI spaces with Carnot tangents, reinforcing the deep structural constraints linking measure-theoretic differentiability, nilpotent approximations, and algebraic topology of groups.  

The statement is simple—“can the Carnot tangent type vary?”—but its resolution would have far-reaching effects across geometric analysis, making it a quintessential example of a “narrow entrance, profound depth” research problem.